\begin{tomtat}
	\begin{minipage}[t]{0.3\textwidth}
		\vspace{0pt}
		\begin{tomtat}
			\begin{center}
				\textbf{ĐỀ 27}
			\end{center}
			
		\end{tomtat}
	\end{minipage}
	\hfill
	\begin{minipage}[t]{0.5\textwidth}
		\vspace{0pt}
		\begin{tomtat}
			\begin{center}
				\includegraphics[scale=0.08]{images/27}
			\end{center}
		\end{tomtat}
	\end{minipage}
\end{tomtat}
\textbf{I. PHẦN TRẮC NGHIỆM}
\begin{ex}
	Giá trị của x thỏa mãn $3(x-1)+2x=7$ là:
	\choice{$x=1$}{$x=0$}{$x=2$}{$x=4$}
	\loigiai{
		Ta có:
		\begin{align*} 3(x-1)+2x &= 7 \\ 3x - 3 + 2x &= 7 \\ 5x - 3 &= 7 \\ 5x &= 7+3 \\ 5x &= 10 \\ x &= \frac{10}{5} \\ x &= 2 \end{align*}
		\\ \textbf{Đáp Án C}
	}
\end{ex}

\begin{ex}
	Cho tam giác ABC vuông tại A, có $AB=12$cm, $BC=13$cm. Tính $\sin C$.
	\choice{$\sin C = \frac{12}{13}$}{$\sin C = \frac{13}{12}$}{$\sin C = \frac{5}{13}$}{$\sin C = \frac{5}{12}$}
	\loigiai{
		Trong tam giác ABC vuông tại A:
		$\sin C = \frac{\text{cạnh đối}}{\text{cạnh huyền}} = \frac{AB}{BC}$.
		Với $AB=12$cm và $BC=13$cm:
		$\sin C = \frac{12}{13}$.
		\\ \textbf{Đáp Án A}
	}
\end{ex}

\begin{ex}
	Thể tích V của hình cầu có bán kính $R=3$cm (tính theo $\pi$) là:
	\choice{$V=36\pi \text{ (cm}^3\text{)}$}{$V=12\pi \text{ (cm}^3\text{)}$}{$V=9\pi \text{ (cm}^3\text{)}$}{$V=27\pi \text{ (cm}^3\text{)}$}
	\loigiai{
		Thể tích của hình cầu được tính bằng công thức $V = \frac{4}{3}\pi R^3$.
		Với $R=3$ cm, ta có:
		\begin{align*} V &= \frac{4}{3}\pi (3)^3 \\ &= \frac{4}{3}\pi \cdot 27 \\ &= 4\pi \cdot 9 \\ &= 36\pi \text{ (cm}^3\text{)} \end{align*}
		\\ \textbf{Đáp Án A}
	}
\end{ex}

\begin{ex}
	Cho đường tròn tâm O. Biết rằng A, B, C là ba điểm phân biệt thuộc đường tròn (O), sao cho $\widehat{BOC}=90^\circ$. Tính số đo $\widehat{BAC}$?
	\choice{$\widehat{BAC}=90^\circ$}{$\widehat{BAC}=45^\circ$}{$\widehat{BAC}=180^\circ$}{$\widehat{BAC}=60^\circ$}
	\loigiai{
		Góc ở tâm $\widehat{BOC}$ và góc nội tiếp $\widehat{BAC}$ cùng chắn cung BC.
		Do đó, số đo góc nội tiếp bằng một nửa số đo góc ở tâm cùng chắn một cung:
		$\widehat{BAC} = \frac{1}{2} \widehat{BOC}$.
		Với $\widehat{BOC}=90^\circ$:
		$\widehat{BAC} = \frac{1}{2} \cdot 90^\circ = 45^\circ$.
		\\ \textbf{Đáp Án B}
	}
\end{ex}

\begin{ex}
	Cho bảng tần số điểm thi kiểm tra giữa kì I môn Toán của 36 học sinh lớp 9A, của một trường THCS.
	\begin{center}
		\begin{tabular}{|c|c|c|c|c|c|c|c|}
			\hline
			Điểm & 4 & 5 & 6 & 7 & 8 & 9 & 10 \\
			\hline
			Tần số & 2 & 3 & 7 & 10 & 9 & 4 & 1 \\
			\hline
		\end{tabular}
	\end{center}
	Tần số tương đối của giá trị điểm 7 là:
	\choice{$27,8$}{$27\%$}{$28\%$}{$27,8\%$}
	\loigiai{
		Tổng số học sinh (tổng tần số) là $N=36$.
		Tần số của giá trị điểm 7 là $n_7 = 10$.
		Tần số tương đối của giá trị điểm 7 được tính bằng công thức:
		$f_7 = \frac{n_7}{N} \times 100\%$.
		$f_7 = \frac{10}{36} \times 100\% = \frac{5}{18} \times 100\%$.
		$\frac{5}{18} \approx 0,2777...$
		$f_7 \approx 0,2777... \times 100\% \approx 27,77...\%$.
		Làm tròn đến một chữ số thập phân, ta được $27,8\%$.
		\\ \textbf{Đáp Án D}
	}
\end{ex}
\begin{ex}
	Mẫu số liệu dưới đây ghi tiền điện (nghìn đồng) của 20 hộ gia đình trong một tháng ở một xóm nhỏ:
	% Dữ liệu: 150; 170; 570; 660; 770; 830; 1100; 660; 160; 470; 620; 380; 400; 720; 590; 950; 640; 159; 780; 700
	Ghép các số liệu thành 4 nhóm ứng với 4 nửa khoảng: [150; 450); [450; 750); [750; 1050); [1050; 1450). Tần số ghép nhóm của nửa khoảng [450; 750)?
	\choice{$8$}{$9$}{$5$}{$10$}
	\loigiai{
		Dữ liệu tiền điện (20 hộ, sắp xếp để dễ đếm):
		150; 159; 160; 170; 380; 400; % Nhóm 1: [150; 450)
		470; 570; 590; 620; 640; 660; 660; 700; 720; % Nhóm 2: [450; 750)
		770; 780; 830; 950; % Nhóm 3: [750; 1050)
		1100 % Nhóm 4: [1050; 1450)
		
		Ta đếm tần số cho mỗi nhóm:
		\begin{itemize}
			\item Nhóm $[150; 450)$: Bao gồm các giá trị $x$ sao cho $150 \le x < 450$.
			Các giá trị: 150; 159; 160; 170; 380; 400.
			Tần số: 6.
			\item Nhóm $[450; 750)$: Bao gồm các giá trị $x$ sao cho $450 \le x < 750$.
			Các giá trị: 470; 570; 590; 620; 640; 660; 660; 700; 720.
			Tần số: 9.
			\item Nhóm $[750; 1050)$: Bao gồm các giá trị $x$ sao cho $750 \le x < 1050$.
			Các giá trị: 770; 780; 830; 950.
			Tần số: 4.
			\item Nhóm $[1050; 1450)$: Bao gồm các giá trị $x$ sao cho $1050 \le x < 1450$.
			Các giá trị: 1100.
			Tần số: 1.
		\end{itemize}
		Tần số ghép nhóm của nửa khoảng $[450; 750)$ là 9.
		\\ \textbf{Đáp Án B}
	}
\end{ex}

\begin{ex}
	Cho hình nón có bán kính đáy bằng 4cm, độ dài đường sinh bằng 40cm. Diện tích xung quanh của hình nón (lấy $\pi \approx 3,14$) là:
	\choice{$251,2 \text{ (cm}^2\text{)}$}{$502,4 \text{ (cm}^2\text{)}$}{$125,6 \text{ (cm}^2\text{)}$}{$1004,8 \text{ (cm}^2\text{)}$}
	\loigiai{
		Bán kính đáy của hình nón là $r = 4$ cm.
		Độ dài đường sinh của hình nón là $l = 40$ cm.
		Diện tích xung quanh của hình nón được tính bằng công thức $S_{xq} = \pi r l$.
		$S_{xq} = \pi \cdot 4 \cdot 40 = 160\pi \text{ (cm}^2\text{)}$.
		Với $\pi \approx 3,14$:
		$S_{xq} \approx 160 \cdot 3,14 = 502,4 \text{ (cm}^2\text{)}$.
		\\ \textbf{Đáp Án B}
	}
\end{ex}

\begin{ex}
	Biết rằng quãng đường s của một vật rơi tự do (bỏ qua sức cản không khí) được tính theo công thức $s = 4,9t^2$, trong đó t là thời gian rơi (s tính bằng mét, t tính bằng giây).
	Thả một vật rơi 10 giây thì vật chạm đất. Độ cao của vật đó từ điểm thả vật đến mặt đất là:
	\choice{$49 \text{ m}$}{$490 \text{ m}$}{$100 \text{ m}$}{$1000 \text{ m}$}
	\loigiai{
		Công thức tính quãng đường rơi tự do là $s = 4,9t^2$.
		Thời gian rơi của vật là $t = 10$ giây.
		Độ cao của vật từ điểm thả vật đến mặt đất chính là quãng đường vật rơi được.
		$s = 4,9 \cdot (10)^2$
		$s = 4,9 \cdot 100$
		$s = 490$ (mét).
		\\ \textbf{Đáp Án B}
	}
\end{ex}
\newpage
\textbf{II. PHẦN TỰ LUẬN}
\begin{ex}
	Cho $a>0$ và $a \neq 1$. Rút gọn biểu thức $P = \left( \frac{1}{\sqrt{a}+1} - \frac{a}{a+\sqrt{a}} \right) : \frac{\sqrt{a}-1}{2\sqrt{a}+2}$.
	\loigiai{
		Điều kiện: $a>0, a \neq 1$.
		Ta có:
		\begin{align*} P &= \left( \frac{1}{\sqrt{a}+1} - \frac{a}{\sqrt{a}(\sqrt{a}+1)} \right) : \frac{\sqrt{a}-1}{2(\sqrt{a}+1)} \\ &= \left( \frac{1}{\sqrt{a}+1} - \frac{\sqrt{a}}{\sqrt{a}+1} \right) \cdot \frac{2(\sqrt{a}+1)}{\sqrt{a}-1} \\ &= \left( \frac{1-\sqrt{a}}{\sqrt{a}+1} \right) \cdot \frac{2(\sqrt{a}+1)}{\sqrt{a}-1} \end{align*}
		Vì $a>0$ và $a \neq 1$, nên $\sqrt{a}+1 \neq 0$ và $\sqrt{a}-1 \neq 0$.
		Ta có $1-\sqrt{a} = -(\sqrt{a}-1)$.
		\begin{align*} P &= \frac{-(\sqrt{a}-1)}{\sqrt{a}+1} \cdot \frac{2(\sqrt{a}+1)}{\sqrt{a}-1} \\ &= \frac{-\cancel{(\sqrt{a}-1)}}{\cancel{(\sqrt{a}+1)}} \cdot \frac{2\cancel{(\sqrt{a}+1)}}{\cancel{(\sqrt{a}-1)}} \\ &= -2 \end{align*}
		Vậy $P = -2$.
	}
\end{ex}

\begin{ex}
	Giải hệ phương trình $\begin{cases} 2x-y=1 \\ x+y=5 \end{cases}$.
	\loigiai{
		Ta có hệ phương trình:
		$$ \begin{cases} 2x-y=1 \quad &(1) \\ x+y=5 \quad &(2) \end{cases} $$
		Cộng vế theo vế của phương trình (1) và (2), ta được:
		\begin{align*} (2x-y) + (x+y) &= 1+5 \\ 3x &= 6 \\ x &= \frac{6}{3} \\ x &= 2 \end{align*}
		Thay $x=2$ vào phương trình (2):
		\begin{align*} 2+y &= 5 \\ y &= 5-2 \\ y &= 3 \end{align*}
		Vậy nghiệm của hệ phương trình là $(x;y) = (2;3)$.
	}
\end{ex}

\begin{ex}
	Biết rằng phương trình bậc hai $x^2+4x-2=0$ có hai nghiệm phân biệt $x_1, x_2$. Không giải phương trình, tính giá trị của biểu thức $T=\frac{1}{x_1}+\frac{1}{x_2}+x_1+x_2$.
	\loigiai{
		Phương trình $x^2+4x-2=0$ có $a=1, b=4, c=-2$.
		$\Delta' = b'^2 - ac = (2)^2 - 1(-2) = 4+2=6 > 0$.
		Phương trình có hai nghiệm phân biệt $x_1, x_2$.
		Theo hệ thức Vi-ét, ta có:
		$$ \begin{cases} x_1 + x_2 = -\frac{b}{a} = -\frac{4}{1} = -4 \\ x_1 x_2 = \frac{c}{a} = \frac{-2}{1} = -2 \end{cases} $$
		Ta biến đổi biểu thức T:
		\begin{align*} T &= \frac{1}{x_1} + \frac{1}{x_2} + x_1+x_2 \\ &= \frac{x_2+x_1}{x_1x_2} + (x_1+x_2) \end{align*}
		Thay giá trị từ hệ thức Vi-ét vào biểu thức T:
		\begin{align*} T &= \frac{-4}{-2} + (-4) \\ &= 2 - 4 \\ &= -2 \end{align*}
		Vậy $T=-2$.
	}
\end{ex}
\begin{ex}
	Một ô tô dự định đi từ thành phố X đến thành phố Y cách nhau 200 km. Sau khi đi được $\frac{1}{2}$ quãng đường với vận tốc dự định, ô tô dừng lại nghỉ 15 phút. Để đến Y đúng giờ, trên quãng đường còn lại, ô tô phải tăng vận tốc thêm 20 km/h so với vận tốc dự định. Tính vận tốc dự định của ô tô.
	\loigiai{
		Gọi vận tốc dự định của ô tô là $x$ (km/h). Điều kiện: $x > 0$.
		Thời gian dự định để đi hết quãng đường 200 km là: $t_{dđ} = \frac{200}{x}$ (giờ).
		
		Nửa quãng đường đầu là: $\frac{1}{2} \cdot 200 = 100$ km.
		Thời gian ô tô đi nửa quãng đường đầu với vận tốc dự định là: $t_1 = \frac{100}{x}$ (giờ).
		
		Thời gian ô tô nghỉ là 15 phút = $\frac{15}{60}$ giờ = $\frac{1}{4}$ giờ.
		
		Nửa quãng đường còn lại là: $200 - 100 = 100$ km.
		Vận tốc ô tô đi trên nửa quãng đường còn lại là: $x + 20$ (km/h).
		Thời gian ô tô đi nửa quãng đường còn lại là: $t_2 = \frac{100}{x+20}$ (giờ).
		
		Tổng thời gian thực tế ô tô đi là: $T = t_1 + \text{thời gian nghỉ} + t_2 = \frac{100}{x} + \frac{1}{4} + \frac{100}{x+20}$ (giờ).
		
		Để đến Y đúng giờ, tổng thời gian thực tế phải bằng thời gian dự định:
		$$ \frac{100}{x} + \frac{1}{4} + \frac{100}{x+20} = \frac{200}{x} $$
		Trừ $\frac{100}{x}$ ở cả hai vế:
		$$ \frac{1}{4} + \frac{100}{x+20} = \frac{200}{x} - \frac{100}{x} $$
		$$ \frac{1}{4} + \frac{100}{x+20} = \frac{100}{x} $$
		Quy đồng mẫu số chung là $4x(x+20)$:
		\begin{align*} x(x+20) + 100 \cdot 4x &= 100 \cdot 4(x+20) \\ x^2 + 20x + 400x &= 400(x+20) \\ x^2 + 420x &= 400x + 8000 \\ x^2 + 420x - 400x - 8000 &= 0 \\ x^2 + 20x - 8000 &= 0 \end{align*}
		Giải phương trình bậc hai $x^2 + 20x - 8000 = 0$:
		Ta có $a=1, b=20, c=-8000$.
		$\Delta = b^2 - 4ac = (20)^2 - 4(1)(-8000) = 400 + 32000 = 32400$.
		$\sqrt{\Delta} = \sqrt{32400} = 180$.
		Phương trình có hai nghiệm phân biệt:
		$x_1 = \frac{-b + \sqrt{\Delta}}{2a} = \frac{-20 + 180}{2 \cdot 1} = \frac{160}{2} = 80$.
		$x_2 = \frac{-b - \sqrt{\Delta}}{2a} = \frac{-20 - 180}{2 \cdot 1} = \frac{-200}{2} = -100$.
		
		Đối chiếu với điều kiện $x > 0$, ta nhận nghiệm $x = 80$.
		Vậy, vận tốc dự định của ô tô là 80 km/h.
	}
\end{ex}
\begin{ex}
	Đội văn nghệ của trường có 2 học sinh nam (gồm An và Bình) và 4 học sinh nữ (gồm Hà, Thu, Hân, Vân). Chọn ra ngẫu nhiên hai bạn hát song ca. Tính xác suất của biến cố H: "2 bạn được chọn có cả nam và nữ".
	\loigiai{
		Học sinh nam: N1 (An); N2 (Bình).
		Học sinh nữ: F1 (Hà); F2 (Thu); F3 (Hân); F4 (Vân).
		Tổng số học sinh: $2+4=6$.
		
		Chọn ngẫu nhiên hai bạn hát song ca. Các cặp có thể chọn (không gian mẫu $\Omega$):
		$\Omega = \{ (\text{N1};\text{N2}); (\text{N1};\text{F1}); (\text{N1};\text{F2}); (\text{N1};\text{F3}); (\text{N1};\text{F4}); (\text{N2};\text{F1}); (\text{N2};\text{F2}); (\text{N2};\text{F3}); (\text{N2};\text{F4}); (\text{F1};\text{F2}); (\text{F1};\text{F3}); (\text{F1};\text{F4}); (\text{F2};\text{F3}); (\text{F2};\text{F4}); (\text{F3};\text{F4}) \}$.
		Số phần tử của không gian mẫu là $n(\Omega) = 15$.
		
		Biến cố H: "2 bạn được chọn có cả nam và nữ".
		Các cặp có cả nam và nữ:
		\begin{itemize}
			\item An (N1) với một bạn nữ: $(\text{N1};\text{F1})$; $(\text{N1};\text{F2})$; $(\text{N1};\text{F3})$; $(\text{N1};\text{F4})$.
			\item Bình (N2) với một bạn nữ: $(\text{N2};\text{F1})$; $(\text{N2};\text{F2})$; $(\text{N2};\text{F3})$; $(\text{N2};\text{F4})$.
		\end{itemize}
		Vậy, $H = \{ (\text{N1};\text{F1}); (\text{N1};\text{F2}); (\text{N1};\text{F3}); (\text{N1};\text{F4}); (\text{N2};\text{F1}); (\text{N2};\text{F2}); (\text{N2};\text{F3}); (\text{N2};\text{F4}) \}$.
		Số phần tử của biến cố H là $n(H) = 2 \times 4 = 8$.
		
		Xác suất của biến cố H là:
		$P(H) = \frac{n(H)}{n(\Omega)} = \frac{8}{15}$.
	}
\end{ex}
\begin{ex}
	Cho đường tròn (O), bán kính R ($R>0$) và dây cung BC cố định. Một điểm A chuyển động trên cung lớn BC sao cho tam giác ABC có ba góc nhọn. Kẻ các đường cao AD, BE của tam giác ABC cắt nhau tại H và BE cắt đường tròn (O) tại F (F khác B).
	\begin{enumerate}
		\item[1)] Chứng minh rằng tứ giác DHEC nội tiếp.
		\item[2)] 
		\begin{enumerate}
			\item[a)] Kẻ đường kính AM của đường tròn (O) và OI vuông góc với BC tại I. Chứng minh rằng I là trung điểm của HM.
			\item[b)] Khi BC cố định, xác định vị trí của A trên đường tròn (O) để $DH \cdot DA$ lớn nhất.
		\end{enumerate}
	\end{enumerate}
\end{ex}


\setcounterref{ex}{0}
